% 依赖项（必须放在主文档导言区，即 \begin{document} 之前）：
% \usepackage{tabularray}
% \UseTblrLibrary{varwidth}
% \usepackage{enumitem}
\begingroup
\nopagebreak[4]
\enlargethispage{1cm}
\footnotesize
\centering

\DefTblrTemplate{conthead-text}{default}{（续表）}
\DefTblrTemplate{contfoot-text}{default}{续下页}

\begin{longtblr}[
  caption = {AI智能体主要安全漏洞汇总},
  label   = {tab:agent_vulns},
]{
  width   = \linewidth,
  colspec = {
    |Q[wd=2cm,l,m]
    |Q[wd=2.5cm,l,m]
    |Q[wd=2cm,l,m]
    |Q[wd=0.8cm,l,m]
    |Q[wd=2cm,l,m]
    |X[1,l,m]|
  },
  cells   = {l,m},
  colsep  = 3pt,
  row{1}  = {font=\bfseries},
  rowhead = 1,
  % 每行内容与上下边框之间各保留少量空间；由该行最高单元格决定整行高度。
  rowsep  = 2.5pt,
  measure = vbox,
  stretch = -1,
}
\hline
\textbf{发现精确日期} & \textbf{漏洞名称/事件} & \textbf{影响产品/框架} & \textbf{CVSS} & \textbf{CVE/编号} & \textbf{漏洞影响简述} \\
\hline
% \caption{AI智能体主要安全漏洞}
% \label{tab:agent_vulns} \\
% \hline
% \textbf{\shortstack{发现精确\\日期}} &
% \textbf{\shortstack{漏洞名称\\事件}} &
% \textbf{\shortstack{影响产品\\框架}} &
% \textbf{\shortstack{CVSS\\评分}} &
% \textbf{\shortstack{CVE\\编号}} &
% \centering\arraybackslash\textbf{漏洞影响简述} \\
% \hline
% \endfirsthead
% % 跨页重复表头
% \hline
% \textbf{\shortstack{发现精确\\日期}} &
% \textbf{\shortstack{漏洞名称\\事件}} &
% \textbf{\shortstack{影响产品\\框架}} &
% \textbf{\shortstack{CVSS\\评分}} &
% \textbf{\shortstack{CVE\\编号}} &
% \centering\arraybackslash\textbf{漏洞影响简述} \\
% \hline
% \endhead

\SetCell[c=6]{c,m} \textbf{Critical（严重，CVSS $\geq$ 9.0）} & & & & & \\
\hline

2026-05-06 & OpenClaw OpenShell沙箱文件写入 TOCTOU / symlink swap~\cite{nvdCVE2026_44112} & OpenClaw $<$ 2026.4.22 & 9.6 & CVE-2026-44112 &
\begin{itemize}[leftmargin=*,topsep=0pt,partopsep=0pt,itemsep=1pt,parsep=0pt]
\item OpenShell沙箱文件系统写入存在time-of-check/time-of-use竞态。
\item 攻击者可通过symlink swap将写入重定向到预期mount root之外。
\item 应升级到2026.4.22或包含修复提交的版本。
\end{itemize}
\\
\hline

2026-04-29 & DocsGPT MCP test绕过导致远程代码执行~\cite{nvdCVE2026_26015} & DocsGPT $\geq$ 0.15.0, $<$ 0.16.0 & 9.8 & CVE-2026-26015 &
\begin{itemize}[leftmargin=*,topsep=0pt,partopsep=0pt,itemsep=1pt,parsep=0pt]
\item 攻击者可构造恶意payload绕过“MCP test”行为。
\item 可在官方DocsGPT网站或本地/公开部署中触发任意远程代码执行。
\item 漏洞已在DocsGPT 0.16.0修复。
\end{itemize}
\\
\hline

2026-04-21 & Flowise MCP stdio命令注入/allowlist绕过~\cite{nvdCVE2026_40933} & Flowise/flowise-components $<$ 3.1.0 & 9.9 & CVE-2026-40933 &
\begin{itemize}[leftmargin=*,topsep=0pt,partopsep=0pt,itemsep=1pt,parsep=0pt]
\item 认证攻击者可在Custom MCP配置中添加stdio MCP server并指定任意命令。
\item 既有npx等允许命令可被组合危险参数绕过输入校验，实现OS命令执行。
\item 漏洞已在Flowise 3.1.0修复。
\end{itemize}
\\
\hline

2026-04-15 & Upsonic MCP task命令注入~\cite{nvdCVE2026_30625} & Upsonic 0.71.6 & 9.8 & CVE-2026-30625 &
\begin{itemize}[leftmargin=*,topsep=0pt,partopsep=0pt,itemsep=1pt,parsep=0pt]
\item MCP server/task 创建功能允许用户定义command和args。
\item npm、npx等allowlist命令可通过参数触发任意OS命令执行。
\item 0.72.0起增加Stdio server可直接执行本机命令的风险提示。
\end{itemize}
\\
\hline

2026-04-07 & Langflow未认证代码注入/RCE~\cite{nvdCVE2025_3248} & Langflow $<$ 1.3.0 & 9.8 & CVE-2025-3248 &
\begin{itemize}[leftmargin=*,topsep=0pt,partopsep=0pt,itemsep=1pt,parsep=0pt]
\item \texttt{/api/v1/validate/code}端点存在代码注入漏洞。
\item 远程未认证攻击者可发送构造HTTP请求执行任意代码。
\item 该漏洞已进入CISA KEV目录，建议升级到1.3.0或更高版本。
\end{itemize}
\\
\hline

2026-04-03 & PraisonAI Gateway未授权WebSocket/info访问~\cite{nvdCVE2026_34952} & PraisonAI Gateway $<$ 4.5.97 & 9.1 & CVE-2026-34952 &
\begin{itemize}[leftmargin=*,topsep=0pt,partopsep=0pt,itemsep=1pt,parsep=0pt]
\item Gateway server 在\texttt{/ws}接受WebSocket连接且\texttt{/info}暴露agent topology，均缺少认证。
\item 任意网络客户端可枚举已注册agents，并向agents及其tool sets发送任意消息。
\item 漏洞已在4.5.97修复，应升级并对网关入口加认证与网络隔离。
\end{itemize}
\\
\hline

2026-04-02 & Azure MCP Server关键功能缺失认证~\cite{nvdCVE2026_32211} & Azure MCP Server / Azure Web Apps hosted service & 9.1 & CVE-2026-32211 &
\begin{itemize}[leftmargin=*,topsep=0pt,partopsep=0pt,itemsep=1pt,parsep=0pt]
\item 关键功能缺少认证，未授权攻击者可经网络访问受影响服务。
\item 可造成信息泄露；Microsoft CNA评分还计入完整性影响。
\item 以Microsoft MSRC / NVD记录为准，CVSS 9.1为Microsoft CNA评分。
\end{itemize}
\\
\hline

2026-01-12 & LibreChat MCP stdio任意命令执行~\cite{nvdCVE2026_22252} & LibreChat $<$ v0.8.2-rc2 & 9.9 & CVE-2026-22252 &
\begin{itemize}[leftmargin=*,topsep=0pt,partopsep=0pt,itemsep=1pt,parsep=0pt]
\item MCP stdio transport接受任意命令且缺少有效校验。
\item 任意认证用户可通过单个API请求在容器内以root执行shell命令。
\item 漏洞已在v0.8.2-rc2修复。
\end{itemize}
\\
\hline

2025-12-23 & LangChain dumps/dumpd序列化注入~\cite{nvdCVE2025_68664} & LangChain $<$ 0.3.81 and $<$ 1.2.5 & 9.3 & CVE-2025-68664 &
\begin{itemize}[leftmargin=*,topsep=0pt,partopsep=0pt,itemsep=1pt,parsep=0pt]
\item \texttt{dumps()}/ \texttt{dumpd()}未转义自由字典中的\texttt{lc} 键结构。
\item 用户可控数据在反序列化时可能被当作LangChain对象处理，而非普通数据。
\item 漏洞已在0.3.81和1.2.5修复，应避免反序列化不可信配置。
\end{itemize}
\\
\hline

2025-12-17 & MCP server git复合漏洞~\cite{nvdCVE2025_68143,nvdCVE2025_68144,nvdCVE2025_68145} & Model Context Protocol Servers/\texttt{mcp-server-git} & 8.8/7.1/9.1 & CVE-2025-68143/4/5 &
\begin{itemize}[leftmargin=*,topsep=0pt,partopsep=0pt,itemsep=1pt,parsep=0pt]
\item \texttt{git\_init} 可在任意可访问路径创建仓库，使后续Git操作越出预期范围。
\item \texttt{git\_diff} / \texttt{git\_checkout} 对用户参数缺少校验，flag-like参数可导致任意文件覆盖。
\item \texttt{--repository} 限制下未校验后续 \texttt{repo\_path} 是否仍在允许仓库内，可操作其他可访问仓库。
\end{itemize}
\\
\hline

2025-09-22 & Flowise CustomMCP 远程代码执行~\cite{nvdCVE2025_59528} & Flowise 3.0.5 & 10.0 & CVE-2025-59528 &
\begin{itemize}[leftmargin=*,topsep=0pt,partopsep=0pt,itemsep=1pt,parsep=0pt]
\item CustomMCP节点解析\texttt{mcpServerConfig}时将用户输入传入。\texttt{Function()}构造器执行
\item 代码在Node.js 运行时权限下执行，可访问\texttt{child\_process}、。\texttt{fs}等危险模块
\item 漏洞已在 Flowise 3.0.6 修复。
\end{itemize}
\\
\hline

2025-07-09 & mcp-remote OAuth authorization endpoint命令注入~\cite{nvdCVE2025_6514} & \texttt{mcp-remote} npm package 0.0.5--0.1.15 & 9.6 & CVE-2025-6514 &
\begin{itemize}[leftmargin=*,topsep=0pt,partopsep=0pt,itemsep=1pt,parsep=0pt]
\item 连接不可信MCP server时，authorization endpoint 响应URL中的构造输入可触发OS命令注入。
\item 攻击需要用户交互但无需攻击者具备权限，成功后影响机密性、完整性和可用性。
\item 已通过\texttt{mcp-remote} 修复提交处置，应升级到不受影响版本。
\end{itemize}
\\
\hline

2025-06-13 & MCP Inspector缺失认证导致远程代码执行~\cite{nvdCVE2025_49596} & MCP Inspector $<$ 0.14.1 & 9.4 & CVE-2025-49596 &
\begin{itemize}[leftmargin=*,topsep=0pt,partopsep=0pt,itemsep=1pt,parsep=0pt]
\item Inspector client 与 proxy之间缺少认证。
\item 未认证请求可经proxy启动MCP stdio命令，造成远程代码执行。
\item 漏洞已在MCPInspector 0.14.1修复。
\end{itemize}
\\
\hline

2025-06-11 & EchoLeak/M365 Copilot AI command injection~\cite{nvdCVE2025_32711} & Microsoft 365 Copilot & 9.3 & CVE-2025-32711 &
\begin{itemize}[leftmargin=*,topsep=0pt,partopsep=0pt,itemsep=1pt,parsep=0pt]
\item M365 Copilot存在AI command injection，未授权攻击者可经网络触发信息泄露。
\item 该漏洞被公开命名为 EchoLeak，NVD 同时引用Microsoft MSRC与Aim Security页面。
\item CVSS 9.3为Microsoft CNA评分；NVD自身评分为7.5。
\end{itemize}
\\
\hline

\SetCell[c=6]{c,m} \textbf{High（高危，7.0 $\leq$ CVSS $\leq$ 8.9）} & & & & & \\
\hline

2026-05-08 & PraisonAI legacy Flask API默认禁用认证~\cite{nvdCVE2026_44338} & PraisonAI 2.5.6 to $<$ 4.6.34 & 7.3 & CVE-2026-44338 &
\begin{itemize}[leftmargin=*,topsep=0pt,partopsep=0pt,itemsep=1pt,parsep=0pt]
\item legacy Flask API server默认禁用认证。
\item 任意可达调用方可访问\texttt{/agents}并通过\texttt{/chat}触发\texttt{agents.yaml}工作流。
\item 漏洞已在4.6.34修复，应升级并限制API暴露面。
\end{itemize}
\\
\hline

2026-04-15 & Agent Zero External MCP Servers远程代码执行~\cite{nvdCVE2026_30624} & Agent Zero 0.9.8 & 8.6 & CVE-2026-30624 &
\begin{itemize}[leftmargin=*,topsep=0pt,partopsep=0pt,itemsep=1pt,parsep=0pt]
\item External MCP Servers配置允许通过JSON指定command和args。
\item 配置应用时这些值会被执行，缺少充分校验或限制。
\item 攻击者可提交恶意MCP配置，以Agent Zero进程权限执行OS命令。
\end{itemize}
\\
\hline

2026-04-15 & Windsurf Prompt Injection to MCP配置命令执行~\cite{nvdCVE2026_30615} & Windsurf 1.9544.26 & 8.0 & CVE-2026-30615 &
\begin{itemize}[leftmargin=*,topsep=0pt,partopsep=0pt,itemsep=1pt,parsep=0pt]
\item 处理攻击者控制的HTML内容时，提示注入可修改本地MCP配置。
\item 可自动注册恶意MCP STDIO server并在无进一步交互下执行任意命令。
\item 成功利用可持久化恶意MC 配置并访问应用暴露的敏感信息。
\end{itemize}
\\
\hline

2026-03-30 & LangChain prompt loading路径遍历~\cite{nvdCVE2026_34070} & LangChain / langchain-core $<$ 1.2.22 & 7.5 & CVE-2026-34070 &
\begin{itemize}[leftmargin=*,topsep=0pt,partopsep=0pt,itemsep=1pt,parsep=0pt]
\item \texttt{langchain\_core.prompts.loading}多个函数未充分校验配置字典中的路径。
\item 当应用把用户可影响的prompt配置传入\texttt{load\_prompt()}或 \texttt{load\_prompt\_from\_config()} 时，攻击者可读取主机文件。
\item 读取受模板/示例文件扩展名检查限制；漏洞已在1.2.22修复。
\end{itemize}
\\
\hline

2026-01-09 & WeKnora MCP stdio 命令注入~\cite{nvdCVE2026_22688} & Tencent WeKnora $<$ 0.2.5 & 8.8 & CVE-2026-22688 &
\begin{itemize}[leftmargin=*,topsep=0pt,partopsep=0pt,itemsep=1pt,parsep=0pt]
\item 认证用户可向MCP stdio设置注入\texttt{stdio\_config.command} / \texttt{args}。
\item 服务端会基于注入值启动subprocess，造成命令执行。
\item 漏洞已在WeKnora 0.2.5修复。
\end{itemize}
\\
\hline

2025-11-07 & Open WebUI Direct Connections SSE代码注入~\cite{nvdCVE2025_64496} & Open WebUI $<$ 0.6.35 & 8.0 & CVE-2025-64496 &
\begin{itemize}[leftmargin=*,topsep=0pt,partopsep=0pt,itemsep=1pt,parsep=0pt]
\item Direct Connections功能允许恶意外部模型服务器通过SSE execute events在受害者浏览器执行JavaScript。
\item 可导致认证token窃取和账户接管；与 Functions API链接时可造成后端RCE。
\item 攻击需受害者启用Direct Connections并添加恶意模型URL；0.6.35已修复。
\end{itemize}
\\
\hline

\SetCell[c=6]{c,m} \textbf{Medium（中危，4.0 $\leq$ CVSS $\leq$ 6.9）} & & & & & \\
\hline

\SetCell[c=6]{c,m} 本表核验条目无中危漏洞。 & & & & & \\
\hline
\end{longtblr}
\endgroup
